% sections/abstract.tex -- the abstract's text (paper_v2, cut in C3 from the full version's).
% B1 (25 November) brought in the author's review of the Vietnamese abstract (paper_v2_vn, L10):
% two ways of choosing the paths with random draws as the control, and a verdict that depends on
% the context. Every number is \val{} from numbers.tex.
HAN and similar models read a heterogeneous graph through meta-paths, usually chosen by hand. Paths
found by the LMSPS search improved HAN on node classification, but whether such a search is worth
its cost in recommendation is untested. On Amazon-Meta, we compare two ways of choosing HAN's
meta-paths from the same \val{search-cands} candidates: by hand (fixed in advance) and by LMSPS
adapted to link prediction (combined from \val{search-runs} runs), with random draws as the
control. Every set is tested in HAN and metapath2vec. The searched sets trail the manual ones by
\val{han-6-hit-relabs}--\val{han-4-hit-relabs}\% on the metric fixed in advance and only tie them on
full ranking, after about \val{search-total-h} GPU hours of search. The loss sits mostly on books
and music, which the searched sets' category paths cannot reach. The random control shows that the
choice of paths moves HAN little on this schema: all sets of a size lie within \val{han-4-span}
Hit@10 of each other, and the random draws do at least as well as both ways. For HAN on this small
schema, automatic search comes fairly close to manual design, so the choice between them depends on
the context: manual design is better when an expert can do it in less time than the search takes,
while search gives nearly equivalent sets without domain knowledge.
